\chapter{Rust Runner：为什么需要本地交易所仿真}

\section{本章要解决的困惑}

如果策略这么简单，为什么还要 Rust Runner？答案是：Runner 不是为了让策略更高级，而是为了让成交更可信。

fast backtest 直接在表格上计算 entry/exit。Runner 则扮演本地交易所：

\begin{lstlisting}
market truth
-> replay clock
-> public stream to Bot
-> order ingress from Bot
-> latency schedule
-> order arrival
-> fill
-> private stream
-> event log
\end{lstlisting}

它解决的是“策略意图在一个可审计交易所里会发生什么”，不是“策略因子怎么设计”。

\section{Rust 新人先理解哪些概念}

不用先学完整 Rust 教材。读本项目的 Runner，只需要几个局部概念：

\begin{longtable}{p{0.24\textwidth}p{0.66\textwidth}}
\toprule
概念 & 本项目里怎么理解 \\
\midrule
\code{struct} & 一组字段，类似 Python dataclass \\
\code{enum} & 多种事件/命令类型中的一种 \\
\code{Result<T, E>} & 可能成功也可能失败的返回 \\
\code{serde} & JSON/NDJSON 序列化和反序列化 \\
\code{clap} & CLI 参数定义 \\
\code{socket} & TCP public/private/order stream \\
\code{stream loop} & 按 replay clock 顺序处理市场和订单 \\
\bottomrule
\end{longtable}

\section{CLI 不是策略}

入口：

\begin{lstlisting}
systems/ccusdt_replay_exchange/engine/crates/cli/src/main.rs
\end{lstlisting}

\code{main.rs} 的主要作用是把命令行参数解析成 Runner 能执行的配置。它告诉系统跑哪种命令：

\begin{lstlisting}
catalog scan
canonical build
run sparse-python
run server
serve
\end{lstlisting}

不要在 CLI 里找策略 alpha。策略在 Python Bot；CLI 只是启动和配置。

\section{Runner 的责任区域}

入口：

\begin{lstlisting}
systems/ccusdt_replay_exchange/engine/crates/runner/src/lib.rs
\end{lstlisting}

不要逐行从头读。按责任区域读：

\begin{longtable}{p{0.28\textwidth}p{0.62\textwidth}}
\toprule
责任 & 你要理解的问题 \\
\midrule
market loading & canonical quote/trade/L2 如何进入 replay \\
public stream & Runner 给 Bot 发送哪些 market events \\
order ingress & Bot 的 submit order 怎么进入 Runner \\
latency queue & observed time 到 arrival time 如何计算 \\
fill model & market buy/sell 用哪个 bid/ask 成交 \\
portfolio & position、cash、equity 如何更新 \\
event log & 哪些事件写进 \code{events.ndjson} \\
\bottomrule
\end{longtable}

\section{top-of-book taker fill}

当前默认成交模型：

\begin{lstlisting}
top_of_book_taker_ioc_v1
\end{lstlisting}

人话：

\begin{itemize}[leftmargin=2em]
  \item market buy 到达时打 arrival ask；
  \item market sell 到达时打 arrival bid；
  \item 默认 \code{fee_bps=0}，但 spread 仍然存在；
  \item maker 字段是保留，不是当前 active maker 模型。
\end{itemize}

\section{observed quote、arrival quote、fill price}

Runner 的价值之一，是把这三个时刻分开：

\begin{longtable}{p{0.24\textwidth}p{0.66\textwidth}}
\toprule
对象 & 含义 \\
\midrule
observed quote & Bot 决策时看到的 quote \\
arrival quote & order 延迟后到达交易所时的 quote \\
fill price & 根据 side 和 arrival quote 选出的成交价 \\
\bottomrule
\end{longtable}

如果 latency 为 0 且 deterministic timer 没问题，observed 和 arrival 应该对齐；如果有延迟或 next-event stress，它们可能不同。

\section{Runner 不读 label}

Runner 只应该知道市场事实、订单、成交、账户。它不能知道这个 entry 未来赚不赚钱。

\warningline{Runner 不读 label。未来收益、MFE、MAE、scored entries 都不能进入 Runner 决策路径。}

\section{手动检查点}

读 strict run 时，找一笔订单：

\begin{lstlisting}
order_intent
order_scheduled
order_arrived
order_ack / order_reject
fill_created
portfolio_state_on_change
\end{lstlisting}

如果这些事件能连起来，你就理解了 Runner 在做什么。

